
{\centering{\capspace{\Huge{\MakeUppercase{introduction}.}}}\par}

\baselineskipvspace

\begin{enpars}
The service of exposition of the Blessed Sacrament regularly takes place in Catholic churches around the world. While the concluding portion of the service known as Benediction, wherein the priest imparts the blessing of Christ the Eucharistic Lord on the faithful assembled in the church, is regulated as to its contents, the rites of exposition, the devotions celebrated immediately before Benediction (the rest of a prolonged period of adoration being passed in silence, as if in the Garden on the evening of the Passion), and the musical choices, both with respect to Gregorian chant and modern hymns, are left to local custom.

Consequently the need arose to gather this material along with a suitable English translation for the parts which are always prayed in Latin according to the rites of the Roman church or which are customarily sung in Latin, namely, the other devotional chants or hymns in honor of the Blessed Virgin Mary and the saints or to mark the liturgical season. Other prayers, in English, are customarily prayed at the Church of the Assumption and needed to be placed alongside the Latin chants and prayers. These are marked by the day of the week or season of the year or month on or in which they are prayed.

Likewise, the chants with their melodies are indicated as assigned to Sundays (that is, all Sundays in the case of prayers; in the case of the \textit{O Salutaris, Tantum ergo,} and the {Adoremus in æternum,} Sundays after Epiphany through Quinquagesima, the last before Ash Wednesday, then Sundays after Pentecost), by the other liturgical seasons (Advent and Lent as well as Paschal Time) or by major feasts (to include the octaves of Christmas and the Epiphany and Pentecost) as well as to weekdays after Epiphany through Quinquagesima, then after Pentecost.

The faithful are invited to sing; certain instructions are given in the book as to which parts are sung by the cantors of the \textit{schola cantorum} or by the faithful. Those who wish to follow in silence may still profit from the song given unto the Lord with the English translations provided. A mix of prose and poetical translations are given for the hymns, according to the merits of each. Some freer translations of prose texts are retained for familiarity. Otherwise, the best prose translation available to use without restrictions was selected for prose texts. But since the English is never to be sung, and only texts given only in English are to be recited, so long as these translations are serviceable, coherent, and as easy to understand as possible, they will be of benefit to all.

The music of Gregorian pieces is largely taken from the Solesmes editions, which have not only the customary melodies but also the customary texts traditionally used at Benediction (e.g. the litanies, the prayers for the pope or the bishop, etc.) except for \textit{Corde natus ex parentis} which retains the melody of \textit{Divinum Mysterium,} a medieval trope on the \textit{Sanctus} of the Mass, a lovely pairing of a medieval melody which is nevertheless not one of the authentic melodies of this hymn; the \textit{Adoremus in æternum} sung in Advent and Lent, taken from the \textit{Westminster Hymnal}; and the hymn \textit{Cælitum Joseph,} which comes from an eighteenth-century French source. Finally, although the \textit{Stabat Mater} and \textit{Adeste Fideles} are given in Gregorian notation, rest assured that their melodies are the familiar ones.
\end{enpars}